\begin{lemma}[Situation 7.3: maximum-index normalization]
Fix a rung-local minimum hierarchy counterexample \(C\), put
\[
 r=\noderef{HierarchyPathNumber}(C.\mathit{hierarchy}),\qquad
 \ell=\noderef{HierarchyFullTailLength}(C),
\]
and assume \(2\le r<\ell\).  Assume also that every vertex has a missing
color under every proper coloring that leaves exactly \(e_0\) uncolored.
Then every \(i\) in
\(\noderef{HierarchyEligibleCommonMissingIndices}\) satisfies \(i<r\).
Thus the maximum eligible common-missing index, when the finite set is
nonempty, is strictly before the path part of the tail.
\end{lemma}
\begin{proof}
Write \(I\) for the finite set in the statement and suppose that
\(i\in I\) with \(r\leq i\).  The nonempty finite set
\[
                   I_{\geq r}=\{t\in I:r\leq t\}
\]
has a greatest member \(m\).  By the definition
\noderef{HierarchyEligibleCommonMissingIndices}, membership of \(m\)
supplies an actual \noderef{HierarchyReplacementPresentation} \(P\) and a
color \(\alpha\), not merely an abstract coloring, such that, for
\(\mu=P.\mathit{coloring}\),
\[
 \alpha\in\noderef{MissingColors}(G,\mu,y_m)\cap
              \noderef{MissingColors}(G,\mu,y_\ell),\qquad
                         r\leq m\leq\ell-1.                    \tag{1}
\]
Let \(\mathcal E_P=P.\mathit{corresponding}.\mathit{replacementETT}\),
let \(H=P.\mathit{counterexample}.\mathit{hierarchy}\), let
\(q=\operatorname{Fin.last}(H.\mathit{dividerCount})\), and put
\[
 A=\noderef{HierarchyStarredStage}(H,q),\qquad
 D=\noderef{HierarchyStoredColorsAt}(H,q),\qquad
 D_n=\noderef{TashkinovStoredColors}(\mathcal E_P,n).
\]
The \noderef{SameGoodHierarchy} field of \(P\), including its underlying
\noderef{SameHierarchyStages} data, identifies the divider count, every
ordered stage, every divider position, all computed stored sets, and every
paper-defined Gamma pair with the corresponding data for \(C\).  Thus the
last starred stage is the same exact tree \(A\), and the definitions of
\noderef{HierarchyFullTailLength}, \noderef{HierarchyFullTailPrefix}, and
\noderef{HierarchyPathNumber} give the same \(\ell\), the same exact
prefixes \(T(y_j)\), and the same path number \(r\) for \(P\).

Assume first that \(m<\ell-1\), and abbreviate
\[
 X^-=T(y_m),\qquad X=T(y_{m+1}),\qquad B=T(y_{\ell-1}).          \tag{2}
\]
The base-stage, nested, strict, core-inside-first, and terminal-stage fields
of \noderef{TashkinovHierarchy}, together with the PE and non-PE clauses of
\noderef{HierarchyBase}, give
\[
 A\prec X^-\prec X\preccurlyeq B\prec T,\qquad |X|=|X^-|+1.    \tag{3}
\]
In the divider-zero PE case the first relation uses the left-segment field
of the \noderef{OrderedTreeUnion} in the parallel base; otherwise it is a
direct stage inclusion.  Every comparison in (3) is an entrywise
\noderef{TreeSequenceInitialSegment}, so it preserves the individual
parallel-edge identities.  Since \(r\leq m\), the candidate clause in
\noderef{HierarchyPathNumber} gives, in particular, the exact consecutive
endpoint pair for the individual edge joining \(y_m\) to \(y_{m+1}\).

Apply \noderef{HierarchyCounterexamplePenultimateElementary} to, in order,
the properness, exactly-one-uncoloredness, complete lower-rung six-pack,
fewer-divider induction, lexicographic-minimality, and \(2n+11\) fields of
the \noderef{HierarchyMinimumCounterexampleContext}
\(P.\mathit{minimumContext}\).  Its conclusion is the derived fact
\[
       \noderef{ColoredVertexSetElementary}
         (G,\mu,\noderef{TreeSequenceVertexSet}(B)).             \tag{4}
\]
The same context says that \(\mu\) is
\noderef{IsProperEdgeColoring} and leaves exactly \(e_0\) uncolored in the
sense of \noderef{IsOneEdgeUncolored}.  Apply the palette-surplus hypothesis
to choose
\[
                 \theta\in\noderef{MissingColors}(G,\mu,y_{m+1}). \tag{5}
\]
The vertices \(y_m,y_{m+1}\) are distinct by the vertex-injectivity field of
\noderef{TreeSequence}.  Hence (1), (4), and (5) give
\(\alpha\ne\theta\), and elementarity says that \(y_m\), respectively
\(y_{m+1}\), is the only vertex of \(B\) missing \(\alpha\), respectively
\(\theta\).  Unfolding the least-candidate definition
\noderef{HierarchyFirstMissingPosition} therefore gives
\[
 T.\mathit{vertex}(i_\alpha)=y_m,\quad
 (i_\alpha).\mathit{val}=|A|+m,\qquad
 T.\mathit{vertex}(i_\theta)=y_{m+1},\quad
 (i_\theta).\mathit{val}=|A|+m+1,
 \qquad
 \alpha,\theta\notin\noderef{TreeMissingColors}(\mu,A).         \tag{6}
\]
At the last stored index, the definition is the literal identity
\[
 D=D_n\setminus\noderef{TreeMissingColors}(\mu,A).              \tag{7}
\]
Consequently nonmembership of either \(\alpha\) or \(\theta\) in \(D\)
implies its nonmembership in \(D_n\).  The theorem
\noderef{HierarchyTailMissingColorExclusion}, applied at the two exact
first positions in (6), excludes each color from its preceding displayed
tail interval.  Its only exception is the individual divider-zero SE/RE
seed when that edge has a color in \(D_n\); equation (7) eliminates this
exception for an unstored affected color.

We record two elementary component facts used at each switch.  Suppose
\(c,d\) have distinct missing witnesses \(u,v\in V(B)\), their full
\((c,d)\)-\noderef{TwoColorComponent}s agree, and \(w\) is a third distinct displayed vertex
missing one of \(c,d\).  By
\noderef{MissingColorIncidenceEndpoint}, each of the three vertex-nodes has
degree at most one in \noderef{TwoColorIncidenceGraph}; for a vertex
missing the second color, interchange the two color arguments.  If all
three lay in one component, \noderef{TwoColorComponentSpanningPath} would
span it by a simple path, while
\noderef{TwoColorIncidenceDegreeAtMostTwo} permits only the two ends of
that path to have degree at most one.  This contradiction, followed by
\noderef{TwoColorComponentsDisjoint}, proves that the component through
\(w\) is disjoint from the common witness component.  The endpoint lemma
also makes the component through \(w\) a possibly singleton
\noderef{TwoColorPathComponent}.

Next suppose \(\psi'\) is obtained from a proper coloring \(\psi\) by
switching such a disjoint terminal component \(Q\), and that \(B\) is
elementary under \(\psi\).  Then, pointwise,
\[
 \noderef{MissingColors}(G,\psi',v)=
 \noderef{MissingColors}(G,\psi,v)\qquad(v\in V(B)).             \tag{8}
\]
Indeed, if \(\operatorname{inl}v\notin Q\), no incident edge-node is in
\(Q\), so every incident value is fixed.  If
\(\operatorname{inl}v\in Q\), disjointness makes \(v\) different from the
two first missing witnesses.  By elementarity, both affected colors are
present at \(v\).  Expanding \noderef{PresentColors}, choose the incident
edge of each color; properness makes each edge identity unique, connectedness
puts both edge-nodes in \(Q\), and \noderef{TwoColorComponentSwap} exchanges
the two values while fixing every other incident color.  This proves (8)
even when \(Q\) meets \(B\) at internal vertices; no vertex-disjointness is
being inferred from component disjointness.

We now give the route construction that will be invoked after every
Cases 1/2 switch.  Let \(R\) be the current replacement presentation, write
\(H_R=R.\mathit{counterexample}.\mathit{hierarchy}\) and
\(\mathcal E_R=R.\mathit{corresponding}.\mathit{replacementETT}\), let
\(\psi\) be its coloring, and let a terminal \((c,d)\)-component switch give
\(\psi'\).  Suppose \noderef{StableChangeCasesOneTwo} has first been
applied to this exact current presentation, so that it returns
\[
 \noderef{HierarchyWeaklyStable}(H_R,q,\psi')\quad\hbox{and}\quad
 \noderef{ModColoring}(\mathcal E_R,\psi'),                    \tag{9}
\]
and suppose the individual edges from \(A\) through \(X\) and the
pointwise equality (8) have been checked for this switch.

Inspect the actual terminal certificate
\[
 \mathcal E_R.\mathit{series}.\mathit{certified}
       (\operatorname{Fin.last}\mathcal E_R.\mathit{rounds})
\]
from \noderef{TashkinovSeries} and
\noderef{TashkinovSeriesStep}.  In PE its
\noderef{TashkinovClosure} contains a \(\psi\)-
\noderef{TashkinovAugmentationSequence} from the terminal ladder tree
\(T_n^R\) to the final ladder tree.  In SE/RE it supplies the exact
individual edge \(f_n^R\), the literal singleton equation
\(F_n^R=\{f_n^R\}\), an exact
\noderef{TreeSequenceAdjoinEdge} seed \(S_R=T_n^R+f_n^R\), and the route
from that seed.  The \noderef{ExtendedTashkinovTree}.belowFinal field puts
\(X\preccurlyeq T\) below the output of this certificate.

The source is below \(A\) in every positive-divider case.  For PE this
follows from the base stage and the nested hierarchy stages.  For SE/RE,
\noderef{HierarchyNonparallelFirstEdge} identifies the certified seed
below the first ordinary stage and hence below \(A\).  At divider zero,
the parallel clause of \noderef{HierarchyBase} and its exact
\noderef{OrderedTreeUnion} put \(T_n^R\preccurlyeq A\) in PE; in SE/RE the
nonparallel clause gives \(A=T_n^R\), and the one-edge seed lies below
\(X\) because \(m\geq r\geq2\).  Thus
\noderef{TashkinovAugmentationSequenceRestrict} restricts the certified
route first to \(T\), then to \(X\), and, whenever the source is below
\(A\), to \(A\).

For clarity, factor rather than merely name the remaining pieces.  Induct
on the reflexive-transitive chain in
\noderef{TashkinovAugmentationSequence}.  Each
\noderef{TashkinovAugmentation} raises the length by exactly one and
retains all preceding entries.  Starting with a source of length at most
\(|A|\leq|X|\) and ending above \(X\), there is a unique chain member of
length \(|A|\) and a unique one of length \(|X|\); comparison as exact
prefixes of the common final tree identifies them with \(A\) and \(X\).
This gives the old-color pieces source-to-\(A\), \(A\)-to-\(X\), and the
actual suffix
\[
                         X\longrightarrow T                    \tag{10}
\]
under \(\psi\).  In divider-zero SE/RE there is no source-to-\(A\) piece;
restrict directly from the exact seed to \(X\), using the reflexive route
if the seed equals \(X\).

At positive divider count, unfold
\noderef{HierarchyWeaklyStable} in (9).  Its positive branch is a
\noderef{TreeColoringInvariant} on the last starred stage \(A\): every
individual displayed nonroot edge of \(A\) is fixed and every vertex
missing set on \(A\) is unchanged.  The root \(e_0\) is fixed because both
colorings leave exactly it uncolored.  Take the exact prefix \(A^-\) of
length \(|A|-1\); positivity of \noderef{TreeSequence} makes it defined,
and invariance restricts to pointwise missing-set equality on \(A^-\).
Hence
\noderef{TashkinovAugmentationSequenceColoringTransport} transports the
source-to-\(A\) piece from \(\psi\) to \(\psi'\).

At divider-zero PE, unfolding the zero branch of
\noderef{HierarchyWeaklyStable} gives exactly
\noderef{CombinedStableColoring} on \(T_n^R\) and the closure \(R_n^R\).
Apply \noderef{CombinedStableClosurePreserved}: it fixes every displayed
edge identity of both \(T_n^R\) and \(R_n^R\), preserves the closure under
\(\psi'\), and its two stable constituents preserve the vertex missing
sets needed by their augmentation routes.  For an arbitrary position
\(s<|A|\), the \noderef{OrderedTreeUnion} from
\noderef{HierarchyBase}.parallel either identifies the edge at \(s\) with
an edge of \(T_n^R\) by its left-segment field, or supplies an exact edge
of \(R_n^R\) by its new-entry-from-right field.  Its corresponding vertex
union gives the same coverage for missing sets.  Thus every edge and
missing-set premise for the source-to-\(A\) transport is proved, not merely
asserted.  In divider-zero SE/RE the exact certified seed is retained and
there is no route below it.

Transport the \(A\)-to-\(X\) piece, or the entire seed-to-\(X\) piece in
the last case, by
\noderef{TashkinovAugmentationSequenceColoringTransport}.  The checked
edge equality on the interval \(|A|\leq s<|X|\) supplies its first
pointwise premise; (8), restricted along \(X^-\preccurlyeq B\), supplies
its missing-set premise on the exact penultimate prefix \(X^-\).  Concatenate
the transported source-to-\(A\) and \(A\)-to-\(X\) pieces using transitivity
of \noderef{TashkinovAugmentationSequence}.  This constructs the required
route
\[
 \begin{cases}
 T_n^R\longrightarrow X,&\text{PE},\\
 T_n^R+f_n^R\longrightarrow X,&\text{SE/RE},
 \end{cases}                                                   \tag{11}
\]
with the exact seed and singleton connecting-edge equation in the second
line.  Finally, when both affected colors are missing on \(X\) after the
switch, \noderef{TashkinovAugmentationSequenceComponentSwap} transports
the old suffix (10) to \(\psi'\).  Concatenating it explicitly with (11)
gives the full route to \(T\).  In particular, the \(hrouteX\) and
\(hsuffix\) inputs of \noderef{StableChangeCasesOneTwoSplitTail} are the
certified pieces (11) and (10), not a claim that all tail edges were fixed.

We also record how the output is returned to the original replacement
family.  The split-tail theorem supplies a corresponding ETT relative to
\(\mathcal E_R\), maximum property, a base, a hierarchy, a
\noderef{GoodHierarchy}, and same-good data.  Compose its
\noderef{CorrespondingExtendedTashkinovTree} with
\(R.\mathit{corresponding}\) field by field: compose rung equalities and,
at every positive ladder index, the dependent series presentations,
ladder-tree equalities, raw stored-set equalities, exact connecting-edge
sets, and extension kinds.  Compose the two positive-stage
\noderef{StableColoring} relations with \noderef{sc2}; the needed protected
sets agree by the just-composed raw data, terminal-ladder missing sets, and
the definition of \noderef{TashkinovStoredColors}.  Compose the two
\noderef{SameGoodHierarchy} witnesses through their
\noderef{SameHierarchyStages}, stored-set, and Gamma-pair fields.

If the switched coloring has one affected color missing at a retained
vertex and at \(y_\ell\), these two distinct witnesses and
\noderef{ColoredVertexSetElementary} make the returned terminal tree
non-elementary, hence give a \noderef{HierarchyCounterexample}.  Its
\noderef{HierarchyMinimumCounterexampleContext} is reconstructed field by
field: properness comes from \noderef{TwoColorComponentSwapProper}, the
swap fixes \(\mathsf{none}\), correspondence transports the complete
lower-rung package and \(2n+11\) estimate, same-good stage equality
transports the fewer-divider assertion, and equality of the relevant path
number and terminal excess identifies
\noderef{HierarchyCounterexampleMeasure}, so the old
lexicographic-minimality field transports.  This gives an actual
original-relative replacement presentation, which we use below after
each switch.

We now treat the stored-color alternatives.

First suppose \(\alpha\notin D\).  If also \(\theta\notin D\), (7) puts
both colors outside \(D_n\), so the divider-zero seed exception is absent.
The tail exclusion for \(\alpha\), including the individual edge into
\(y_{m+1}\) because that edge is incident with the
\(\alpha\)-missing vertex \(y_m\), supplies the exact unused interval from
\(i_\alpha\) to \(i_\theta\).  The relevant alternative of
\noderef{ChangeLemma} identifies the full \((\alpha,\theta)\)-components
through \(y_m\) and \(y_{m+1}\).  The three-endpoint argument above
separates that component from the terminal component through \(y_\ell\).
Switch the latter component.  Equations (3), (6), (7), and the exact unused
interval are precisely the external-first-missing hypotheses of
\noderef{StableChangeCaseThree}; that theorem itself returns the full
corresponding ETT, maximum property, rebuilt good hierarchy, and same-good
data.  Package its output by the fieldwise construction above.  The switch
makes \(\theta\) missing at \(y_\ell\), or leaves it already missing there
in the singleton case, and separation preserves (5).  Thus
\(m+1,P',\theta\) witnesses \(m+1\in I\).

Assume next that \(\alpha\notin D\) but \(\theta\in D\).  By the
\(\mathit{twoColors}\) field of \noderef{GoodHierarchy}, choose
\(\gamma\in P.\mathit{counterexample}.\mathit{good}.\Gamma^q_\theta\).
Its avoidance field, read
through \noderef{HierarchySegmentColors}, puts \(\gamma\) in
\noderef{HierarchyStageMissingColors}(H,q)\) and excludes it from every
exact edge from \(A\) through the first \(\theta\)-missing position
\(y_{m+1}\).  At divider zero that stage-missing witness is on \(T_n\),
which is contained in \(A\) by the parallel ordered union or equals \(A\)
in the non-PE case.  Thus \(\gamma\) is missing on \(A\), so (7) gives
\(\gamma\notin D\); conversely \(\alpha\), being nonmissing on \(A\), lies
in no last Gamma pair.

Apply \noderef{ChangeLemma} to \((\gamma,\alpha)\).  It identifies the
component through the first \(\gamma\)-missing stage vertex with the one
through \(y_m\).  Separate the terminal component, switch it, call the new
coloring \(\mu_1\), and immediately apply
\noderef{StableChangeCasesOneTwo}.  Avoidance fixes every possible
\(\gamma\)-edge from \(A\) through \(X\); tail exclusion and
\(\alpha\notin D_n\) fix every preceding \(\alpha\)-edge, while
\(\alpha\)-missingness at \(y_m\) fixes the sole next edge into \(X\).
This proves
the individual fixed-prefix premise.  The inside/outside argument (8)
proves pointwise missing-set equality on \(B\), and both affected colors
remain missing on \(X\) at their common first-witness component.

Here are the two universal split-tail premises for this first switch.
Fix arbitrary \(\eta\in D\).  The affected colors \(\gamma,\alpha\) are
both outside \(D\), so \(\eta\) differs from both.  Their component swap
does not change the missing predicate for \(\eta\) at any vertex.  Unfold
\noderef{HierarchyFirstMissingPosition}: if the candidate set is nonempty,
the two identical predicates give identical candidate finite sets and the
same least member; if it is empty, both colorings use the identical last
position as the totalized fallback.  Hence
\[
 \operatorname{First}_{\mu_1}(\eta)=\operatorname{First}_{\mu}(\eta)
                         \qquad(\eta\in D).                     \tag{12}
\]
For the Gamma cutoff, \(\alpha\) belongs to no pair.  If
\(\gamma\in\Gamma^q_\eta\), it also belongs to
\(\Gamma^q_\theta\); the disjoint-pairs field of
\noderef{GoodHierarchy} forces \(\eta=\theta\).  Its old first missing
position is exactly the last vertex of \(X\), so it is at most \(|X|\).
This proves the cutoff for every stored \(\eta\).  Apply the certified
route construction (9)--(11), and then
\noderef{StableChangeCasesOneTwoSplitTail} with (12) and this cutoff.  Its
output packages to an original-relative presentation \(P_1\), under whose
coloring \(\gamma\) is missing both at its retained stage witness and at
\(y_\ell\).

Reapply \noderef{HierarchyCounterexamplePenultimateElementary} to all six
fields of \(P_1.\mathit{minimumContext}\); \(B\) is elementary for the
intermediate coloring.  The same-good data preserve the exact trees,
stored set \(D\), and pair \(\Gamma^q_\theta\).  Write \(D^1\) for the
last stored set of \(P_1\); the same-good equality says \(D^1=D\).
Under \(\mu_1\), apply
\noderef{ChangeLemma} to \((\gamma,\theta)\), identify the stage-witness
component with the component through \(y_{m+1}\), separate the component
through \(y_\ell\), switch it, and apply
\noderef{StableChangeCasesOneTwo} before using weak stability.  New
avoidance excludes \(\gamma\) through \(X\).  Tail exclusion excludes
\(\theta\) through its exact first position \(y_{m+1}\).  If the
divider-zero SE/RE seed has color \(\theta\in D_n^1\), its actual
\noderef{TashkinovSeriesStep} makes it incident with \(T_n^1\): this is
the inside/other-endpoint data of \noderef{MaximumDefectiveChoice} in SE
and of \noderef{RevisitingExtensionCertificate} in RE.  It is therefore a
\noderef{TreeIncidentEdge}, and the protected-edge clause of the returned
\noderef{StableColoring} fixes it.  Thus the fixed-prefix premise holds
also in this exceptional seed case, and (8) gives the new pointwise
missing-set equality on \(B\).

For the second stored-position premise, fix arbitrary
\(\eta\in D^1\).  Stage missingness excludes \(\gamma\) from \(D^1\).  If
\(\eta\ne\theta\), it is unaffected, and the nonempty/empty analysis of
\noderef{HierarchyFirstMissingPosition} used in (12) applies verbatim.  If
\(\eta=\theta\), pointwise equality on \(B\) preserves the actual witness
\(y_{m+1}\); intermediate elementarity says no earlier vertex is a witness,
so the least position remains \(|A|+m+1\), rather than falling into the
fallback branch.  Thus every stored first position agrees.  The affected
color \(\theta\) is not stage-missing and hence belongs to no Gamma pair;
if \(\gamma\in\Gamma^q_\eta\), pairwise disjointness with the retained
\(\Gamma^q_\theta\) again gives \(\eta=\theta\), whose first position is
the last vertex of \(X\).  This is the second universal Gamma cutoff.
Invoke the route construction for the actual intermediate terminal
certificate, concatenate its transported pieces, and apply the split-tail
theorem.  The resulting original-relative \(P_2\) has \(\theta\) missing
at both \(y_{m+1}\) and \(y_\ell\), so \(m+1\in I\).

It remains in the propagation case that \(\alpha\in D\).  Let
\(\Gamma^q_\alpha=\{\gamma_1,\gamma_2\}\), with
\(\gamma_1\ne\gamma_2\), using \noderef{GoodHierarchy}.twoColors.  The
individual edge \(e_{m+1}\) joining \(y_m\) to \(y_{m+1}\) has at most one
color.  Relabel the two reserved colors so that
\[
                       \mu(e_{m+1})\ne\operatorname{some}\gamma_1. \tag{13}
\]
This is the one-edge choice made in the paper.  Avoidance for
\(\Gamma^q_\alpha\) excludes \(\gamma_1\) from every preceding edge through
\(y_m\), and (13) extends that exclusion across the exact next edge and
hence through \(X\).  Tail exclusion for \(\alpha\), plus its missingness
at the incident vertex \(y_m\), excludes \(\alpha\) through \(X\), except
possibly the named divider-zero SE/RE seed.

Apply \noderef{ChangeLemma} to \((\gamma_1,\alpha)\), separate its common
stage/\(y_m\) component from the terminal component, switch the latter,
and call the coloring \(\nu_1\).  Apply
\noderef{StableChangeCasesOneTwo} immediately.  Equation (13) and Gamma
avoidance fix every \(\gamma_1\)-candidate edge through \(X\).  If the sole
possible seed has stored color \(\alpha\), the actual SE or RE certificate
again makes it a \noderef{TreeIncidentEdge} of \(T_n\), so the
\noderef{StableColoring} edge clause fixes it.  Thus every exact edge from
\(A\) through \(X\) is fixed.  Equation (8) supplies missing-set equality
on \(B\), and both \(\gamma_1,\alpha\) remain missing on \(X\).

For \(h_{\rm storedFirst}\), take arbitrary \(\eta\in D\).  Stage
missingness gives \(\gamma_1\notin D\).  If \(\eta\ne\alpha\), the color is
unaffected and the candidate set and fallback argument from (12) applies.
If \(\eta=\alpha\), (8) preserves its actual first witness \(y_m\), and
elementarity excludes any earlier witness, so its first position is still
\(|A|+m\).  For \(h_{\rm gammaCutoff}\), \(\alpha\) itself belongs to no Gamma
pair because (6) says it is not stage-missing.  If
\(\gamma_1\in\Gamma^q_\eta\), disjointness with
\(\Gamma^q_\alpha\) forces \(\eta=\alpha\), whose first position
\(|A|+m\) is
strictly before the last vertex of \(X\).  These are separate quantified
proofs of the two premises.  Use (9)--(11) for this exact switch and apply
\noderef{StableChangeCasesOneTwoSplitTail}.  Package the result as \(P_1\);
\(\gamma_1\) is missing at its retained stage witness and at \(y_\ell\),
whereas \(\alpha\) remains missing at \(y_m\).

Under \(P_1\), rederive elementarity of \(B\) from all six fields of its
minimum context.  Its same-good data preserve the pair
\(\Gamma^q_\alpha\), the avoidance interval through \(y_m\), and (13),
because the first switch fixes \(e_{m+1}\).  The color \(\theta\) is still
missing first at \(y_{m+1}\).  Again write \(D^1\) for its last stored set;
same-good equality identifies it with \(D\).  Apply \noderef{ChangeLemma} to
\((\gamma_1,\theta)\): Gamma avoidance together with the explicit
one-edge exclusion (13), now under \(\nu_1\), is its complete unused
interval through \(y_{m+1}\).  Separate and switch the terminal component,
then apply \noderef{StableChangeCasesOneTwo}.

For this second switch, the same avoidance plus (13) fixes every
\(\gamma_1\)-candidate edge through \(X\), and tail exclusion fixes every
\(\theta\)-candidate edge.  If \(\theta\) is a stored divider-zero seed
color, use the actual SE/RE certificate and the protected incident-edge
clause exactly as above.  Equation (8), now with the intermediate
elementary coloring, gives pointwise equality on \(B\).  Fix arbitrary
\(\eta\in D^1\).  The color \(\gamma_1\) is not stored.  If
\(\eta\ne\theta\), its missing predicate is globally unaffected; if
\(\eta=\theta\), equality on \(B\) preserves its actual first witness
\(y_{m+1}\) and excludes an earlier one.  Unfolding
\noderef{HierarchyFirstMissingPosition}, including the identical fallback
for every unaffected no-witness color, proves all stored-position
equalities.  The color \(\theta\) is not stage-missing and lies in no Gamma
pair.  If \(\gamma_1\in\Gamma^q_\eta\), disjointness with the copied
\(\Gamma^q_\alpha\) forces \(\eta=\alpha\); its first witness \(y_m\) is in
\(X\).  This proves the universal cutoff.  Construct the actual
intermediate route by (9)--(11), transport and concatenate its suffix, and
apply the split-tail theorem.  The resulting original-relative \(P_2\)
has \(\theta\) missing at \(y_{m+1}\) and \(y_\ell\).

In every branch for \(m<\ell-1\), we have therefore constructed an actual
replacement presentation witnessing \(m+1\in I\).  Since \(m+1>m\) and
\(m+1\geq r\), this contradicts the definition of \(m\) as the greatest
member of \(I_{\geq r}\).

It remains to exclude \(m=\ell-1\).  Keep \(P,\mu,H,A\) as above and put
\[
 B=T(y_{\ell-1}),\qquad X^-=T(y_{\ell-2}),\qquad
 e_\ell=T.\mathit{edge}(|T|-1).                                \tag{14}
\]
Here \(2\leq r\leq\ell-1\), so \(\ell\geq3\).  The candidate equation for
\(r\) in \noderef{HierarchyPathNumber} gives the exact endpoint identity
\[
                 G.\mathit{ends}(e_\ell)=\{y_{\ell-1},y_\ell\}. \tag{15}
\]
Independently reapply
\noderef{HierarchyCounterexamplePenultimateElementary} to the properness,
exactly-one-uncoloredness, complete lower-rung package, fewer-divider
induction, lexicographic-minimality, and large-missing-set fields of
\(P.\mathit{minimumContext}\).  Thus \(V(B)\) is elementary under \(\mu\)
also in this branch; the properness and exactly-one-uncoloredness used
below are the first two fields of that same
\noderef{HierarchyMinimumCounterexampleContext}.

Inspect the same actual terminal \noderef{TashkinovSeriesStep}.  Its PE
\noderef{TashkinovClosure} gives a \(\mu\)-augmentation route from \(T_n\);
its SE/RE constructors give the exact individual seed \(T_n+f_n\), its
singleton connecting-edge equation, and the route from that seed.  The
upper-containment field of \noderef{ExtendedTashkinovTree} puts \(T\),
\(B\), and \(X^-\) below the final closure.  The hierarchy inclusions and,
in non-PE, exact seed length put the applicable source below \(B\).
Restrict with \noderef{TashkinovAugmentationSequenceRestrict} first to
\(T\) and then to \(B\).

Peel the last step from the finite source-to-\(T\) chain.  Every
\noderef{TashkinovAugmentation} increases length by one, and
\(|T|=|B|+1\); hence the final predecessor has length \(|B|\).  Both it and
\(B\) are exact prefixes of \(T\), so they agree entry by entry.  Unfolding
that last augmentation supplies a color \(\beta\) with
\[
 \mu(e_\ell)=\operatorname{some}\beta,\qquad
 \beta\in\noderef{TreeMissingColors}(\mu,B).                   \tag{16}
\]
Choose a displayed \(z\in V(B)\) missing \(\beta\).  The edge \(e_\ell\)
is incident with \(y_{\ell-1}\) by (15) and has color \(\beta\) by (16), so
\noderef{MissingColors} gives \(z\ne y_{\ell-1}\).  The exact one-entry
relation \(X^-\prec B\), together with vertex injectivity in
\noderef{TreeSequence}, therefore puts \(z\in V(X^-)\).

Define
\[
 \sigma(e)=
 \begin{cases}
   \operatorname{some}\alpha,&e=e_\ell,\\
   \mu(e),&e\ne e_\ell.
 \end{cases}                                                   \tag{17}
\]
The color \(\alpha\) is missing under \(\mu\) at both endpoints in (15).
For two \noderef{AdjacentEdges}, whose common endpoint is supplied by
\noderef{Incident}, properness is inherited if neither is
\(e_\ell\); if one is \(e_\ell\), its common endpoint is one of those two
vertices and missingness excludes \(\alpha\) from the other edge.  Hence
\(\sigma\) is \noderef{IsProperEdgeColoring}.  The displayed index of
\(e_\ell\) is positive, so edge injectivity and the first-edge field of
\noderef{TreeSequence} give \(e_\ell\ne e_0\).  Equation (16) and the fact
that (17) fixes \(\mathsf{none}\) show directly from
\noderef{IsOneEdgeUncolored} that exactly \(e_0\) remains uncolored.

Neither endpoint in (15) belongs to \(V(X^-)\), and no displayed edge of
\(B\) is \(e_\ell\).  Therefore
\[
 \begin{aligned}
 \sigma(B.\mathit{edge}(t))&=\mu(B.\mathit{edge}(t))
                                      &&(t<|B|),\\
 \noderef{MissingColors}(G,\sigma,v)&=
       \noderef{MissingColors}(G,\mu,v)&& (v\in V(X^-)).
 \end{aligned}                                                 \tag{18}
\]
In particular, every edge incident with the terminal ladder tree \(T_n\)
is fixed and every vertex missing set on \(T_n\) agrees.  Thus \(\sigma\)
is a \noderef{StableColoring} of \(T_n\) with protected set \(D_n\).

Apply \noderef{GoodHierarchyTerminalPrefixRepackage} first, under the old
coloring \(\mu\), to
\[
 (\mathcal E_P,P.\mathit{maximum},P.\mathit{replacementBase},H,
   P.\mathit{counterexample}.\mathit{good},B).
\]
The strict hypothesis \(A\prec B\) follows from \(\ell\geq3\), and
\(B\preccurlyeq T\) is exact.  Denote its outputs by
\(\mathcal E_B,\mathcal B_B,H_B,g_B\).  The theorem states that
\(\mathcal E_B\) retains the complete certified series and rung of
\(\mathcal E_P\), has maximum property, that \(\mathcal B_B\) retains the
old core, closure, and numerical support, and that \(H_B,g_B\) retain every
nonterminal stage, divider, computed stored set, stage-missing set, and
Gamma pair.  Only the terminal tree is shortened to \(B\).

Transport the already restricted old route to \(B\) with
\noderef{TashkinovAugmentationSequenceColoringTransport}, taking \(X^-\)
as the exact penultimate prefix.  The edge premise is the first line of
(18), and the vertexwise missing-set premise is its second line.  We obtain
\[
 \begin{cases}
 T_n\longrightarrow B,&\text{PE},\\
 T_n+f_n\longrightarrow B,&\text{SE/RE},
 \end{cases}                                                   \tag{19}
\]
under \(\sigma\), retaining in the second case the actual seed and its
singleton equation.

We now build the new ETT and base, with a separate certified construction
in the PE case.  Suppose the terminal kind is PE.  The core of
\(\mathcal B_B\) lies in \(A\preccurlyeq X^-\), and its parallel
\noderef{OrderedTreeUnion} says that the vertex set of the core is the
union of those of \(T_n\) and the retained closure \(R_n\).  Hence every
vertex of \(T_n\) and \(R_n\), and every edge incident with either tree,
is disjoint from the changed edge \(e_\ell\).  Equation (18) therefore
verifies both \noderef{StableColoring} constituents of
\[
 \noderef{CombinedStableColoring}(T_n,R_n,D_n,\mu,\sigma).      \tag{20}
\]
By \noderef{CombinedStableClosurePreserved}, the same ordered \(R_n\) is
a \(\sigma\)-closure of the same support prefix, and every displayed edge
of \(T_n\) and \(R_n\) is fixed.  Apply the Lean-closed theorem
\noderef{HierarchySplitter} to the old-color prefix package
\((\mathcal E_B,\mathcal B_B,H_B)\), its maximum property and positive
rung, the lower-rung field of \(P.\mathit{minimumContext}\), properness and
one-uncoloredness of \(\mu,\sigma\), the PE identity, (20), and the actual
route in the first line of (19).  It returns one actual
\noderef{CorrespondingExtendedTashkinovTree} on \(B\), maximum property
for that same replacement ETT, a \(\sigma\)-base with the retained core and
closure, and a hierarchy \(K\) satisfying
\[
                  \noderef{SameHierarchyStages}(H_B,K).        \tag{21}
\]
Thus no equality of independently chosen PE support witnesses is assumed:
the comparison of the two singleton connecting edges and the inside/outside
endpoint fields of \noderef{MaximumDefectiveChoice} is exactly the content
already proved in \noderef{HierarchySplitter}.

Suppose instead that the terminal kind is SE or RE.  Apply \noderef{Td} to
\(\mathcal E_B\), its maximum property and positive rung, the complete
lower-rung field of \(P.\mathit{minimumContext}\), and the just-proved
terminal stability of \(\sigma\).  This gives a
\noderef{ModColoring}.  Instantiate its universal clause with the exact
seed, singleton equation, and second route in (19), obtaining a
corresponding ETT on \(B\); then
\noderef{MaximumPropertyTransport} gives its maximum property.  Define its
base core to be its terminal ladder tree.  Correspondence identifies that
tree with \(T_n\), while \noderef{HierarchyBase}.nonparallel identifies the
old core with \(T_n\).  Copy the unused closure and numerical support;
terminal missing-color compatibility is reflexive under \(\sigma\), the PE
implication is vacuous, and the non-PE implication is the defining core
equality.  Copy the divider count, ordered stages, and positions of \(H_B\).
All \noderef{TashkinovHierarchy} fields are exact-prefix, strict-length, or
entrywise equations, so this constructs \(K\) and the same-stage relation
(21) also in the non-PE case.

It remains to build goodness on \(K\).  Correspondence identifies the raw
stored arrays.  Equation (18) identifies missing sets on \(T_n\), every
ordinary stage, and every starred stage, since they lie in
\(A\preccurlyeq X^-\).  By the definitions of
\noderef{TashkinovStoredColors}, \noderef{HierarchyStoredColorsAt}, and
\noderef{HierarchyStageMissingColors}, the corresponding computed sets for
\(H_B\) and \(K\) are equal.  Copy every Gamma pair through (21).

We check the only nonautomatic bound, the totalized first missing position
on the shortened terminal tree \(B\).  For a color other than
\(\alpha,\beta\), (17) changes no incident edge of that color, so its
missing predicate is identical at every vertex and the nonempty least
candidate, or the empty fallback, is identical.  For \(\beta\), the
unchanged witness \(z\in X^-\) occurs strictly before the only vertex whose
missingness can change, so its least missing position is unchanged.  For
\(\alpha\), (4) says its old first witness in \(B\) is the last vertex
\(y_{\ell-1}\).  Under \(\sigma\) the new \(\alpha\)-edge makes it present
there and there is no earlier witness; hence the no-witness branch of
\noderef{HierarchyFirstMissingPosition} returns that same last numerical
position.  Thus every color has the same totalized first position under
\(\mu\) on \(H_B\) and under \(\sigma\) on \(K\).  Together with the fixed
edge identities of \(B\), this gives equality of every
\noderef{HierarchySegmentColors} set.

The two-color and disjoint-pair fields of \noderef{GoodHierarchy} now copy.
Avoidance copies from the equal stage-missing and segment-color sets; every
unique-refresh field copies from the equal stages, stored sets, and Gamma
pairs.  Every positive-stage \noderef{TreeColorSetMinusClosed} assertion
lies inside \(X^-\), so (18) fixes its missing sets and boundary-edge
values.  In the PE parallel-base field, the core and closure also lie in
\(X^-\); (20) and \noderef{CombinedStableClosurePreserved} transport its
non-minus-closed assertion and cardinality bound.  These checks construct a
genuine good hierarchy \(g_K\) on \(K\).

Finally, \(z\) still misses \(\beta\) under \(\sigma\).  At
\(y_{\ell-1}\), properness of \(\mu\) makes \(e_\ell\) the unique incident
\(\beta\)-edge, and (17) removes its \(\beta\)-value, so that vertex also
misses \(\beta\) under \(\sigma\).  The two vertices are distinct; hence
the new terminal tree \(B\) is not
\noderef{ColoredVertexSetElementary}.  With the positive rung, maximum
property, \(K\), and \(g_K\), this gives a genuine
\noderef{HierarchyCounterexample} \(C'\).

The corresponding certificate gives the same rung, and (21) gives the same
divider count.  The old candidate \(r\) remains a candidate for \(B\): all
its required consecutive endpoint equations except the removed last one
are unchanged, and the removed last one already held by the old candidate.
Conversely, if \(s<r\) were a candidate for \(B\), all its equations in
\(B\) would hold, and the one additional equation for \(e_\ell\) follows
from the old \(r\)-candidate because \(s<r\).  Thus \(s\) would be a
candidate for \(T\), contrary to the minimality in the definition of \(r\).
The new path number is therefore exactly \(r\).

The last ordinary hierarchy stage is unchanged and is strictly below
\(B\).  Since \(|B|=|T|-1\), direct natural-number subtraction gives equal
first coordinates and a second coordinate smaller by one.  By
\noderef{HierarchyCounterexampleLess},
\[
 \noderef{HierarchyCounterexampleMeasure}(C')
   <\noderef{HierarchyCounterexampleMeasure}
          (P.\mathit{counterexample}).                         \tag{22}
\]
This contradicts the same-rung, equal-divider lexicographic-minimality
field of \(P.\mathit{minimumContext}\).  Both possibilities
\(m<\ell-1\) and \(m=\ell-1\) are impossible.  Hence no member of \(I\) is
at least \(r\), which proves the claim.
\end{proof}
